\SourcePage{433}
\SGAsection{12}{关于 $G_{\mathrm{af}}$ 与“反仿射”群的补充}
\footnote{编者注（135）：本节为新增。}我们先给出下述\Term{lemma}{引理}，它把 11.18.1 推广到不假定 $G$ 为\Term{finite-type}{有限型}的情形。\footnote{编者注（136）：本\Term{lemma}{引理}由 M. Raynaud 告知，将用于命题 12.9 的证明。}

\begin{SGAresult}{\Term{lemma}{引理}}{12.1}
设 $k$ 为\Term{field}{域}，$u:G\to H$ 为 $k$-\Term{group}{群}的\Term{monomorphism}{单态射}，其中 $H$ 为\Term{affine}{仿射}。假定 $u$ \Term{quasi-compact}{拟紧}。则 $u$ 是\Term{closed-immersion}{闭浸入}。
\end{SGAresult}
\begin{proof}
利用 (fpqc) \Term{descent}{下降}，可假定 $k$ 为\Term{algebraically-closed}{代数闭域}。根据第 VI$_{\mathrm A}$ 讲 6.4，$u$ 的\Term{closed-image}{闭像} $I$ 是 $H$ 的\Term{closed-group-subscheme}{闭群子概形}，因而仍为\Term{affine}{仿射}。所以用 $I$ 代替 $H$ 后，可假定 $u$ \Term{schematically-dominant}{概形支配}。由于 $k$ 为\Term{algebraically-closed}{代数闭域}，$H'=H_{\mathrm{r\text{é}d}}$ 是 $H$ 的\Term{group-subscheme}{群子概形}；记 $G'=G\times_H H'$，则由 $u$ 经\Term{base-change}{基变换}得到的\Term{morphism}{态射} $u':G'\to H'$ 是\Term{quasi-compact}{拟紧}\Term{monomorphism}{单态射}，并且是\Term{dominant}{支配}的（$u$ 与 $u'$ 的底层\Term{continuous-map}{连续映射}相同）。故根据第 VI$_{\mathrm A}$ 讲 6.2，$u'$ \Term{faithfully-flat}{忠实平坦}；它因而是\Term{faithfully-flat}{忠实平坦}且\Term{quasi-compact}{拟紧}的\Term{monomorphism}{单态射}，从而是\Term{isomorphism}{同构}（参见第 IV 讲 1.14）。

所以 $u:G\to H$ 是\Term{homeomorphism}{同胚}，由 2.9.1，它是\Term{affine}{仿射}的。因此 $G$ 为\Term{affine}{仿射}，再由 11.18.2 (iii)，$u$ 是\Term{closed-immersion}{闭浸入}。
\end{proof}

\begin{SGAresult}{\Term{theorem}{定理}}{12.2}
设 $G$ 为\Term{algebraic-group}{代数 $k$-群}。用 $\rho$ 表示\Term{canonical-morphism}{典范态射} $G\to G_{\mathrm{af}}$，用 $N$ 表示其\Term{kernel}{核}。

(i) \Term{canonical-morphism}{典范态射} $G/N\to G_{\mathrm{af}}$ 是\Term{isomorphism}{同构}，因而 $G_{\mathrm{af}}$ 是\Term{affine}{仿射}\Term{algebraic-group}{代数群}，且 $\rho$ \Term{faithfully-flat}{忠实平坦}。

(ii) 有\Term{canonical-isomorphism}{典范同构} $(G/N)_{\mathrm{af}}=G_{\mathrm{af}}$。

(iii) $N$ 是 $G$ 的\Term{characteristic-subgroup}{特征子群}。

(iv) $\mathcal O(N)=k$。

(v) $N$ \Term{smooth}{光滑}、\Term{connected}{连通}且\Term{commutative-group}{交换}。
\end{SGAresult}
\SourcePage{434}
\begin{proof}
(i) 是 11.18.1 的特例，而 (ii) 由 $G_{\mathrm{af}}$ 与 $(G/N)_{\mathrm{af}}$ 的\Term{universal-property}{泛性质}得出。

证明 (iii)。对任意 $k$-\Term{scheme}{概形} $\pi:S\to\operatorname{Spec}k$，考虑\Term{cartesian-square}{笛卡儿方块}：
\[
\begin{tikzcd}
G_S \arrow[r] \arrow[d,"q"'] & G \arrow[d,"p"]\\
S \arrow[r,"\pi"'] & \operatorname{Spec}k .
\end{tikzcd}
\]
由于 $p$ \Term{quasi-compact}{拟紧}且\Term{separated}{分离}，而 $\pi$ \Term{flat}{平坦}，故根据 EGA III 1.4.15 与 EGA IV$_1$ 1.7.21，有
\[
q_*(\mathcal O_{G_S})=\pi^*(\mathcal O(G))=\pi^*(\mathcal O(G_{\mathrm{af}})),
\]
从而由 EGA II 1.5.2 得 $(G_S)_{\mathrm{af}}=(G_{\mathrm{af}})_S$。因此 $N_S$ 作为\Term{canonical-morphism}{典范态射} $G_S\to(G_S)_{\mathrm{af}}$ 的\Term{kernel}{核}，在 $G_S$ 的任一\Term{automorphism}{自同构}下\Term{invariant}{不变}，即 $N$ 是 $G$ 的\Term{characteristic-subgroup}{特征子群}。

为证明 (iv)，置 $N'=\operatorname{Ker}(N\to N_{\mathrm{af}})$；由 (ii)，它是 $G$ 的\Term{normal-subgroup}{正规子群}。因为 $N$ 是\Term{algebraic-group}{代数群}（它是 $G$ 的\Term{closed-subgroup}{闭子群}），故由 (i)，$N/N'\simeq N_{\mathrm{af}}$；此外，根据第 VI$_{\mathrm A}$ 讲 3.2 及 5.3.2，有 $k$-\Term{group}{群}的\Term{isomorphism}{同构}
\[
(G/N')/N_{\mathrm{af}}\simeq(G/N')/(N/N')\simeq G/N.
\]
由于 $N_{\mathrm{af}}$ 为\Term{affine}{仿射}，根据 9.2 (vii)，\Term{projection}{投影} $G/N'\to G/N$ 也是\Term{affine}{仿射}的；而 $G/N=G_{\mathrm{af}}$ 为\Term{affine}{仿射}，故 $G/N'$ 也为\Term{affine}{仿射}。所以由 $G_{\mathrm{af}}$ 的\Term{universal-property}{泛性质}，\Term{projection}{投影} $p':G\to G/N'$ 经 $G_{\mathrm{af}}=G/N$ \Term{resolution}{分解}，因而 $N\subset N'$，故 $N=N'$。于是 $N_{\mathrm{af}}$ 为\Term{trivial-group}{平凡群}，从而 $\mathcal O(N)=k$。

最后，断言 (v) 是下述\Term{lemma}{引理}的推论。
\end{proof}

\begin{SGAresult}{\Term{lemma}{引理}}{12.3}
设 $k$ 为\Term{field}{域}，$N$ 为满足 $\mathcal O(N)=k$ 的\Term{algebraic-group}{代数 $k$-群}。则 $N$ \Term{smooth}{光滑}、\Term{connected}{连通}且\Term{commutative-group}{交换}。
\end{SGAresult}
\begin{proof}
事实上，可假定 $k$ 为\Term{algebraically-closed}{代数闭域}。此时 $H=N^0_{\mathrm{r\text{é}d}}$ 是 $N$ 的 $k$-\Term{subgroup}{子群}，并且根据第 VI$_{\mathrm A}$ 讲 5.5.1 及 5.6.1，\Term{quotient-scheme}{商 $k$-概形} $X=N/H$ 在 $k$ 上\Term{finite}{有限}（因而\Term{affine}{仿射}）。另一方面，由于 $p:N\to X$ \Term{faithfully-flat}{忠实平坦}，有 $\mathcal O(X)\subset\mathcal O(N)=k$。由此得到 $N=N^0_{\mathrm{r\text{é}d}}$，所以 $N$ \Term{smooth}{光滑}（第 VI$_{\mathrm A}$ 讲 1.3.1）且\Term{connected}{连通}。

再令 $Z$ 为 $N$ 的\Term{center}{中心}。根据 6.2.6，$N/Z$ 为\Term{affine}{仿射}，与上面一样可得 $\mathcal O(N/Z)=k$，从而 $N=Z$。这证明了 12.3，也完成了 12.2 的证明。
\end{proof}

我们还陈述下述\Term{theorem}{定理}而不加证明。（回顾，\Term{field}{域} $k$ 上的\Term{abelian-variety}{阿贝尔簇}是\Term{proper}{固有}、\Term{smooth}{光滑}且\Term{connected}{连通}的 $k$-\Term{group-scheme}{群概形}。）

\begin{SGAresult}{\Term{theorem}{定理}}{12.4（Chevalley）}
设 $k$ 为\Term{perfect-field}{完美域}，$G$ 为\Term{smooth}{光滑}\Term{connected}{连通}的\Term{algebraic-group}{代数 $k$-群}。则存在\Term{affine}{仿射}、\Term{smooth}{光滑}且\Term{connected}{连通}的 $k$-\Term{subgroup}{子群} $L$，它在 $G$ 中\Term{normal-subgroup}{正规}，并且商 $G/L$ 为\Term{abelian-variety}{阿贝尔簇}。此外，$L$ 唯一，其构造与\Term{base-field-extension}{基域扩张}可交换。
\end{SGAresult}

\textbf{注记 12.5。}（1）此\Term{theorem}{定理}由 C. Chevalley 于 1953 年宣布，他在 1960 年发表了证明（[Ch60]）。在此期间，
\SourcePage{435}
I. Barsotti 与 M. Rosenlicht 各自独立地给出了其他证明（[Ba55, Ro56]）；历史方面的评论见 [Se99]。

（2）B. Conrad 给出了 Chevalley 证明的现代版本，即使用\Term{scheme}{概形}语言的版本（[Co02]）。（注意，在该文中，“algebraic group” 指\Term{smooth}{光滑}\Term{connected}{连通}的 $k$-\Term{group-scheme}{群概形}。）

（3）另一方面，Ngô B.-C. 在 2005—2006 年于 Orsay 开设的课程中给出了 Rosenlicht 证明的现代版本。

（4）去掉 $k$ 为\Term{perfect-field}{完美域}的假设后，仍存在最小的\Term{affine}{仿射}\Term{connected}{连通}\Term{normal-subgroup}{正规子群} $L$（未必\Term{smooth}{光滑}），使 $G/L$ 为\Term{abelian-variety}{阿贝尔簇}（[BLR]，§9.2，\Term{theorem}{定理} 1）。

（5）也可去掉 $G$ 在 $k$ 上\Term{smooth}{光滑}的假设：事实上，根据第 VII$_{\mathrm A}$ 讲 8.3，存在整数 $n\geq1$ 使商 $G'=G/({}_{\operatorname{Fr}^n}G)$ \Term{smooth}{光滑}；于是 $G'$ 包含一个如上述（4）中的\Term{subgroup}{子群} $L'$，而 $L'$ 在 $G$ 中的\Term{inverse-image}{逆像}仍具有相同性质。因此，对\Term{field}{域} $k$ 上任一\Term{connected}{连通}\Term{algebraic-group}{代数群} $G$，都存在\Term{exact-sequence}{正合列}
\[
1\longrightarrow H\longrightarrow G\longrightarrow A\longrightarrow1,
\]
其中 $H$ 为\Term{affine}{仿射} $k$-\Term{group}{群}，$A$ 为 $k$-\Term{abelian-variety}{阿贝尔簇}。此外，由 [Per76]，推论 4.2.9，对任一\Term{connected}{连通} $k$-\Term{group}{群} $G$（不必为\Term{algebraic-group}{代数群}）都有这样的\Term{exact-sequence}{正合列}。

（6）设 $k$ 为\Term{algebraically-closed}{代数闭域}，$G$ 为\Term{elliptic-curve}{椭圆曲线} $E$ 与\Term{constant-group}{常 $k$-群} $\{\pm1\}_k$ 对\Term{action}{作用} $(-1)\cdot x=-x$ 所成的\Term{semidirect-product}{半直积}；此时，若 $L$ 是 $G$ 的\Term{closed-subgroup}{闭}\Term{normal-subgroup}{正规子群}且 $G/L$ \Term{connected}{连通}，则 $L=G$。

\textbf{注记 12.6。}依照 [Br09]，若 $\mathcal O(N)=k$，便称 $k$-\Term{group}{群} $N$ 为\Term{anti-affine}{反仿射}（anti-affine）。根据 12.3 与 12.4，若 $k$ 为\Term{perfect-field}{完美域}，则任一\Term{anti-affine}{反仿射}\Term{algebraic-group}{代数 $k$-群}，都是一个\Term{abelian-variety}{阿贝尔簇}被一个\Term{affine}{仿射}、\Term{smooth}{光滑}、\Term{connected}{连通}且\Term{commutative-group}{交换}的\Term{algebraic-group}{代数 $k$-群}的\Term{extension}{扩张}。关于\Term{perfect-field}{完美域}上\Term{anti-affine}{反仿射}\Term{algebraic-group}{代数群}的精确结构及其各种推论，见 M. Brion 与 C.、F. Sancho de Salas 最近的文章（[Br09, SS09]）。

本节最后将证明 M. Raynaud 的两个结果，第一个是注记 11.11.1，第二个是第 XVII 讲附录 III 的命题 2.1。我们需要下述\Term{lemma}{引理}\footnote{编者注（137）：这是 O. Gabber 对 M. Raynaud 告知的一个命题所作的改进版本。}，它改进了 [BAC]，§III.5 中给出的\Term{flatness}{平坦性}判别准则（对于\Term{complete-discrete-valuation-ring}{完备离散赋值环} $R$）。

\begin{SGAresult}{\Term{lemma}{引理}}{12.7}
设 $R$ 为\Term{discrete-valuation-ring}{离散赋值环}，$K$ 为其\Term{fraction-field}{分式域}，$\pi$ 为\Term{uniformizer}{一致化元}。设 $A$ 为\Term{flat}{平坦} $R$-\Term{algebra}{代数}，$M$ 为在 $R$ 上\Term{flat}{平坦}的 $A$-\Term{module}{模}。假定：

(i) $M/\pi M$ 是 $\overline A=A/\pi A$ 上的\Term{flat}{平坦}\Term{module}{模}；

(ii) $M\otimes_R K$ 是 $A\otimes_R K$ 上的\Term{flat}{平坦}\Term{module}{模}。

则 $M$ 是\Term{flat}{平坦} $A$-\Term{module}{模}。
\end{SGAresult}

\SourcePage{436}
\begin{proof}
根据\Term{nilpotent}{幂零}情形的\Term{flatness}{平坦性}判别准则（参见 [BAC]，III §5.2，\Term{theorem}{定理} 1），对每个 $n\in\mathbb N^*$，$M/\pi^nM$ 都是 $A/\pi^nA$ 上的\Term{flat}{平坦}\Term{module}{模}。于是由 [RG71]，II \Term{lemma}{引理} 1.4.2.1，$M$ 是\Term{flat}{平坦} $A$-\Term{module}{模}。为方便读者，简述其证明。置 $s=\pi^n$，$P=M/sM$。由于 $P$ 在 $A/sA$ 上\Term{flat}{平坦}，而后者在 $A$ 上的\Term{projective-dimension}{投射维数}为 $1$，故由\Term{composite-functor}{复合函子}的\Term{spectral-sequence}{谱序列}，$P$ 在 $A$ 上的\Term{tor-dimension}{Tor 维数}不超过 $1$。又由于 $M$ 在 $R$ 上\Term{flat}{平坦}，因而无 $\pi$-\Term{torsion}{挠}，$M_K/M$ 是各 $A$-\Term{module}{模} $\pi^{-n}M/M\simeq M/\pi^nM$ 的\Term{inductive-limit}{归纳极限}，故 $M_K/M$ 的\Term{tor-dimension}{Tor 维数}也不超过 $1$。由\Term{exact-sequence}{正合列} $0\to M\to M_K\to M_K/M\to0$，以及假设 $M_K$ 在 $A_K$ 上从而在 $A$ 上\Term{flat}{平坦}，得到 $M$ \Term{flat}{平坦}。

为完整起见，也给出 O. Gabber 指出的下述较简单证明。设 $I$ 为 $A$ 的\Term{finite-type}{有限生成}\Term{ideal}{理想}，需要证明\Term{morphism}{态射} $u:M\otimes_A I\to M$ 为\Term{injective}{单射}。由假设 (ii)，$u\otimes_RK$ 为\Term{injective}{单射}，所以 $\operatorname{Ker}(u)$ 是一个 $\pi$-\Term{torsion}{挠} $R$-\Term{module}{模}。因此，只需证明 $M\otimes_A I$ 的 $\pi$-\Term{torsion}{挠}部分为零；而它是 $\operatorname{Tor}_1^A(M,I/\pi I)$ 的一个商，这可由\Term{exact-sequence}{正合列}
\[
0\longrightarrow I\xrightarrow{\pi}I\longrightarrow I/\pi I\longrightarrow0
\]
与 $M$ 作\Term{tensor-product}{张量积}看出。

另一方面，由于 $M$ 无 $\pi$-\Term{torsion}{挠}（因为在 $R$ 上\Term{flat}{平坦}），对一切 $i\geq1$，有 $\operatorname{Tor}_i^A(M,\overline A)=0$。所以，若 $(P_\bullet)$ 是 $A$-\Term{module}{模} $M$ 的\Term{projective-resolution}{投射分解}，则 $(P_\bullet\otimes_A\overline A)$ 是 $\overline A$-\Term{module}{模} $\overline M=M/\pi M$ 的\Term{projective-resolution}{投射分解}。因而对任一 $\overline A$-\Term{module}{模} $N$，有
\[
\operatorname{Tor}_i^A(M,N)=\operatorname{Tor}_i^{\overline A}(\overline M,N),
\]
而当 $i\geq1$ 时该式为零，因为 $\overline M$ 在 $\overline A$ 上\Term{flat}{平坦}。故 $\operatorname{Tor}_1^A(M,I/\pi I)=0$，\Term{lemma}{引理}得证。
\end{proof}

\textbf{注记 12.8。}设 $S$ 为\Term{dimension}{维数} $1$ 的\Term{locally-noetherian}{局部诺特}\Term{regular}{正则}\Term{scheme}{概形}，$X$ 为在 $S$ 上\Term{flat}{平坦}、\Term{quasi-separated}{拟分离}且\Term{quasi-compact}{拟紧}的 $S$-\Term{scheme}{概形}。则 $\mathscr A(X)$ 是\Term{flat}{平坦} $\mathcal O_S$-\Term{module}{模}。事实上，可假定 $S$ 为\Term{local-scheme}{局部概形}，记 $s$ 为其\Term{closed-point}{闭点}，$i$ 为包含映射 $X_s\hookrightarrow X$，$\pi$ 为 $R=\mathcal O(S)$ 的\Term{uniformizer}{一致化元}；由于 $X$ 在 $S$ 上\Term{flat}{平坦}，有\Term{sheaf}{层}的\Term{exact-sequence}{正合列}
\[
0\longrightarrow\mathcal O_X\xrightarrow{\pi}\mathcal O_X\longrightarrow i_*(\mathcal O_{X_s})\longrightarrow0,
\]
取\Term{global-sections}{整体截面}即得 $\mathscr A(X)$ 是无 $\pi$-\Term{torsion}{挠}的 $R$-\Term{module}{模}，因而\Term{flat}{平坦}。\footnote{编者注（138）：更一般地，若 $S$ 是\Term{dimension}{维数}不超过 $2$ 的\Term{locally-noetherian}{局部诺特}\Term{regular}{正则}\Term{scheme}{概形}，此结论仍成立，参见 [Ray70a]，VII 3.2。}还可得到，由\Term{morphism}{态射} $X\to X_{\mathrm{af}}$ 诱导的从 $\mathscr A((X_{\mathrm{af}})_s)=\mathscr A(X)/\pi\mathscr A(X)$ 到 $\mathscr A(X_s)$ 的\Term{morphism}{态射}为\Term{injective}{单射}。

现在可以证明下述命题（参见注记 11.11.1）。

\begin{SGAresult}{命题}{12.9}
设 $S$ 为\Term{dimension}{维数}不超过 $1$ 的\Term{locally-noetherian}{局部诺特}\Term{regular}{正则}\Term{scheme}{概形}，$G$ 为在 $S$ 上\Term{flat}{平坦}、\Term{quasi-separated}{拟分离}且\Term{quasi-compact}{拟紧}的 $S$-\Term{group-scheme}{群概形}。则下述条件\Term{equivalence}{等价}：

(i) $G$ 在 $S$ 上\Term{affine}{仿射}。

(ii) $G$ 在 $S$ 上\Term{quasi-affine}{拟仿射}。
\SourcePage{437}
(iii) $G$ \Term{faithful-action}{忠实作用}于一个\Term{quasi-affine}{拟仿射}且\Term{flat}{平坦}的 $S$-\Term{scheme}{概形} $X$。

(iv) $G$ \Term{linear-action}{线性}且\Term{faithful-action}{忠实地作用}于一个在 $S$ 上\Term{flat}{平坦}的\Term{quasi-coherent}{拟凝聚}\Term{module}{模} $\mathcal E$。

(v) \Term{morphism}{态射} $\rho:G\to G_{\mathrm{af}}$ 是\Term{monomorphism}{单态射}。
\end{SGAresult}
\begin{proof}
蕴含 (i) $\Rightarrow$ (ii) 显然，(ii) $\Rightarrow$ (iii) 也显然（取 $X=G$）。

假定 (iii) 成立。由于 $\mathscr A(G)$ 与 $\mathscr A(X)$ 都是\Term{flat}{平坦} $\mathcal O_S$-\Term{module}{模}，按 11.11 的方法可得 $G$ 在 $X_{\mathrm{af}}$ 上\Term{faithful-action}{忠实地右作用}，因而在\Term{flat}{平坦}\Term{quasi-coherent}{拟凝聚} $\mathcal O_S$-\Term{module}{模} $\mathscr A(X)$ 上\Term{linear-action}{线性}且\Term{faithful-action}{忠实地左作用}。

另一方面，若 (iv) 成立，则根据 11.6.1 (ii)，\Term{monomorphism}{单态射} $G\to\underline{\operatorname{Aut}}_{\mathcal O_S}(\mathcal E)$ 经 $G_{\mathrm{af}}$ \Term{resolution}{分解}，故 $G\to G_{\mathrm{af}}$ 是\Term{monomorphism}{单态射}。

最后，假定 (v) 成立，证明 $\rho:G\to G_{\mathrm{af}}$ 是\Term{isomorphism}{同构}。用 $S$ 的一个\Term{connected-component}{连通分支}代替 $S$ 后，可假定 $S$ \Term{irreducible}{不可约}，其\Term{generic-point}{泛点}为 $\eta$。由于 $G_{\mathrm{af}}$ 的构造与\Term{flat}{平坦}\Term{base-change}{基变换}可交换，有 $(G_{\mathrm{af}})_\eta=(G_\eta)_{\mathrm{af}}$，故\Term{morphism}{态射} $G_\eta\to(G_\eta)_{\mathrm{af}}$ 是\Term{monomorphism}{单态射}，根据 12.1 是\Term{closed-immersion}{闭浸入}，又因 $\mathcal O((G_\eta)_{\mathrm{af}})=\mathcal O(G_\eta)$，它是\Term{isomorphism}{同构}。若 $S=\operatorname{Spec}(\kappa(\eta))$，证明已经完成；因此可假定 $\dim S=1$。

取 $S$ 的一个\Term{closed-point}{闭点} $s$，我们证明 $G_s\to(G_{\mathrm{af}})_s$ 是\Term{isomorphism}{同构}，且 $\rho$ 在 $G_s$ 的所有点处\Term{flat}{平坦}。为此，可假定 $S$ 为\Term{local-scheme}{局部概形}，其\Term{closed-point}{闭点}为 $s$。经\Term{base-change}{基变换}得到的\Term{morphism}{态射} $\rho_s:G_s\to(G_{\mathrm{af}})_s$ 是\Term{monomorphism}{单态射}，因而根据 12.1 是\Term{closed-immersion}{闭浸入}；所以由 $\rho_s$ 诱导的\Term{morphism}{态射} $\mathcal O((G_{\mathrm{af}})_s)\to\mathcal O(G_s)$ 为\Term{surjective}{满射}。而根据前一注记，它也是\Term{injective}{单射}，故为\Term{isomorphism}{同构}。（特别地，$\rho$ 因而为\Term{surjective}{满射}。）

于是由\Term{lemma}{引理} 12.7，$\rho:G\to G_{\mathrm{af}}$ \Term{faithfully-flat}{忠实平坦}。由于 $G$ 在 $S$ 上\Term{quasi-compact}{拟紧}，而 $G_{\mathrm{af}}$ 在 $S$ 上\Term{separated}{分离}，故 $\rho$ 也\Term{quasi-compact}{拟紧}（参见 EGA I 6.6.4）。因此 $\rho$ 是\Term{faithfully-flat}{忠实平坦}且\Term{quasi-compact}{拟紧}的\Term{monomorphism}{单态射}，从而是\Term{isomorphism}{同构}。命题得证。
\end{proof}

最后，证明第 XVII 讲附录 III 的命题 2.1；利用 [Per76]，推论 4.2.5，我们把假设中的“\Term{finite-type}{有限型}”替换为“\Term{quasi-compact}{拟紧}且\Term{quasi-separated}{拟分离}”（若假定 $G$ 为\Term{finite-type}{有限型}，便可用 12.1 代替该文的结果）。

\begin{SGAresult}{命题}{12.10}
设 $S$ 为\Term{dimension}{维数}不超过 $1$ 的\Term{locally-noetherian}{局部诺特}\Term{regular}{正则}\Term{scheme}{概形}，$G$ 为\Term{flat}{平坦}、\Term{quasi-compact}{拟紧}且\Term{quasi-separated}{拟分离}的 $S$-\Term{group-scheme}{群概形}。

(i) \Term{canonical-morphism}{典范态射} $\rho:G\to G_{\mathrm{af}}$ \Term{faithfully-flat}{忠实平坦}且\Term{quasi-compact}{拟紧}。因此，$G_{\mathrm{af}}$ \Term{representable}{表示} $G$ 对 $N=\operatorname{Ker}(\rho)$ 的 (fpqc) \Term{quotient-sheaf}{商层}。

(ii) 若 $G$ 还在 $S$ 上为\Term{finite-type}{有限型}，则 $\rho$ 为\Term{finite-presentation}{有限表示}，且 $G_{\mathrm{af}}$ \Term{representable}{表示} $G$ 对 $N$ 的 (fppf) \Term{quotient-sheaf}{商层}，并在 $S$ 上为\Term{finite-type}{有限型}。

(iii) 再假定对 $S$ 的每个\Term{maximal-point}{极大点} $\eta$，$G_\eta$ 都\Term{affine}{仿射}。则 $N$ 是\Term{etale}{平展} $S$-\Term{group}{群}；若 $G$ 在 $S$ 上\Term{separated}{分离}，则 $N$ 是\Term{unit-group}{单位群}。
\end{SGAresult}
\begin{proof}
首先，由于 $G_{\mathrm{af}}$ 在 $S$ 上\Term{affine}{仿射}从而\Term{separated}{分离}，故 $\rho$ \Term{quasi-compact}{拟紧}（参见 EGA I 6.6.4），且\Term{kernel}{核} $N=\operatorname{Ker}(\rho)$ 是 $G$ 的\Term{closed-subgroup}{闭子群}。
\SourcePage{438}
此外，用 $S$ 的一个\Term{connected-component}{连通分支}代替 $S$ 后，可假定 $S$ \Term{irreducible}{不可约}，其\Term{generic-point}{泛点}为 $\eta$。

注意，为证明 (i) 与 (ii)，只需证明 $\rho$ \Term{faithfully-flat}{忠实平坦}；因为这样一来，根据第 IV 讲 5.1.7.1，$G_{\mathrm{af}}$ \Term{representable}{表示} $G$ 对 $N$ 的 (fpqc) \Term{quotient-sheaf}{商层}；若 $G$ 还在 $S$ 上为\Term{finite-type}{有限型}，从而为\Term{finite-presentation}{有限表示}（因为 $S$ \Term{locally-noetherian}{局部诺特}），则由 9.2 (xiii)，$\rho$ 为\Term{finite-presentation}{有限表示}（$G_{\mathrm{af}}\to S$ 亦然），所以 $\rho$ 是 (fppf) \Term{topology}{拓扑}的\Term{covering-morphism}{覆盖态射}。

因此可假定 $S=\operatorname{Spec}(R)$，其中 $R$ 为\Term{discrete-valuation-ring}{离散赋值环}（当 $\dim S=1$），或为\Term{field}{域} $\kappa(\eta)$（当 $\dim S=0$）。记 $s$ 为 $S$ 的\Term{closed-point}{闭点}。由于 $G_{\mathrm{af}}$ 的构造与\Term{flat}{平坦}\Term{base-change}{基变换}可交换，\Term{canonical-morphism}{典范态射} $G_\eta\to(G_\eta)_{\mathrm{af}}$ 可等同于\Term{morphism}{态射} $\rho_\eta:G_\eta\to(G_{\mathrm{af}})_\eta$；又由于
\[
\mathcal O((G_{\mathrm{af}})_\eta)=\mathcal O((G_\eta)_{\mathrm{af}})=\mathcal O(G_\eta),
\]
故根据 [Per76]，4.2.5（也见第 VI$_{\mathrm A}$ 讲的增补 6.6），$\rho_\eta$ \Term{faithfully-flat}{忠实平坦}。若 $\dim S=1$，类似可得 $\rho_s$ \Term{faithfully-flat}{忠实平坦}，因为由注记 12.8，\Term{morphism}{态射} $\mathcal O((G_{\mathrm{af}})_s)\to\mathcal O(G_s)$ 为\Term{injective}{单射}。所以由\Term{lemma}{引理} 12.7，$\rho$ \Term{faithfully-flat}{忠实平坦}。这证明了 (i) 与 (ii)。特别地，$N$ 在 $S$ 上\Term{flat}{平坦}。

现假定 $G_\eta$ 为\Term{affine}{仿射}。由于 $\rho_\eta$ 就是\Term{canonical-morphism}{典范态射} $G_\eta\to(G_\eta)_{\mathrm{af}}$，其\Term{kernel}{核} $N_\eta$ 是\Term{unit-group}{单位群}。我们证明 $N$ 在 $S$ 上\Term{etale}{平展}。由于 $N$ 在 $S$ 上\Term{flat}{平坦}，只需证明对 $S$ 的任一点 $s\ne\eta$，$N_s$ 在 $\kappa(s)$ 上\Term{etale}{平展}。若 $S=\operatorname{Spec}(\kappa(\eta))$，已无须证明，故可假定 $S=\operatorname{Spec}(R)$，其中 $R$ 为\Term{discrete-valuation-ring}{离散赋值环}。设 $s$ 为 $S$ 的\Term{closed-point}{闭点}，$K$ 为 $R$ 的\Term{fraction-field}{分式域}，$\pi$ 为\Term{uniformizer}{一致化元}。取 $x\in N_s$，取 $x$ 在 $N$ 中的一个\Term{affine}{仿射}\Term{open-neighborhood}{开邻域} $U_x$；由于 $N$ 在 $S$ 上\Term{flat}{平坦}，$U_x\cap N_\eta$ \Term{nonempty}{非空}，因而等于 $\{\varepsilon(\eta)\}$，其中 $\varepsilon$ 表示\Term{unit-section}{单位截面}。故 $A=\mathcal O(U_x)$ 是\Term{flat}{平坦} $R$-\Term{algebra}{代数}，满足 $A_K=K$ 且 $\pi^{-1}\notin A$（因为 $\pi$ 属于 $\mathcal O_{U,x}$ 的\Term{maximal-ideal}{极大理想}）。由此得 $A=R$，所以\Term{projection}{投影} $U_x\to S$ 是\Term{isomorphism}{同构}。这证明 $N$ 在 $S$ 上\Term{etale}{平展}；若 $G$ 还在 $S$ 上\Term{separated}{分离}，则逆\Term{isomorphism}{同构} $S\to U_x$ 等于\Term{unit-section}{单位截面}（因为二者在 $S=\operatorname{Spec}(R)$ 的\Term{dense-open}{稠密开集} $\{\eta\}$ 上相同），所以 $N$ 是\Term{unit-group}{单位群}。命题得证。
\end{proof}

特别地，得到下述推论；其另外两个证明见 [An73]，命题 2.3.1 及 [PY06]，命题 3.1。

\begin{SGAresult}{推论}{12.10.1}
设 $R$ 为\Term{discrete-valuation-ring}{离散赋值环}，$K$ 为其\Term{fraction-field}{分式域}，$G$ 为在 $R$ 上\Term{separated}{分离}、\Term{flat}{平坦}且\Term{finite-type}{有限型}的 $R$-\Term{group-scheme}{群概形}。若 $G_K$ 为\Term{affine}{仿射}，则 $G$ 为\Term{affine}{仿射}。
\end{SGAresult}

\textbf{注记 12.10.2。}（a）一方面，O. Gabber 向我们指出一些例子，其中 $G$ 是一个\Term{discrete-valuation-ring}{离散赋值环}上的\Term{flat}{平坦}\Term{finite-type}{有限型}\Term{group}{群}，其\Term{generic-fiber}{泛纤维}是\Term{abelian-variety}{阿贝尔簇}，而 $G\to G_{\mathrm{af}}$ 的\Term{kernel}{核} $N$ 并不\Term{smooth}{光滑}。

（b）另一方面，指出 M. Raynaud 给出了如下例子：取 $S$ 为\Term{field}{域} $k$ 上的 $2$ 维\Term{affine-space}{仿射空间}，有一个\Term{smooth}{光滑}且\Term{quasi-affine}{拟仿射}的 $S$-\Term{group-scheme}{群概形}，其\Term{fiber}{纤维}\Term{affine}{仿射}且\Term{connected}{连通}，但它并不在 $S$ 上\Term{affine}{仿射}；参见 [Ray70a]，§VII.3，第 116 页。

\SourcePage{439}
\SGAsection{13}{维数不超过 $2$ 的正则基底上的平坦仿射群}
\footnote{编者注（139）：本节为新增。}首先指出，证明\Term{field}{域} $k$ 上任一\Term{affine}{仿射}\Term{algebraic-group}{代数群}均为\Term{linear-group}{线性群}的熟知论证，以及\Term{lemma}{引理} 11.12，都可推广到如下情形：$G$ 是在\Term{dimension}{维数}不超过 $2$ 的\Term{noetherian}{诺特}\Term{regular}{正则}\Term{base-scheme}{基概形} $S$ 上\Term{affine}{仿射}、\Term{flat}{平坦}且\Term{finite-type}{有限型}的\Term{group-scheme}{群概形}。当 $S$ 的\Term{dimension}{维数}不超过 $1$ 时，这就证明了注记 11.11.1 的 (b)。推广至 $\dim S=2$ 的情形，要用到 O. Gabber 告知的下述\Term{lemma}{引理}。

\begin{SGAresult}{\Term{lemma}{引理}}{13.1}
设 $S$ 为\Term{noetherian}{诺特}\Term{normal}{正规}\Term{scheme}{概形}，$\mathcal A$ 为\Term{flat}{平坦}\Term{quasi-coherent}{拟凝聚} $\mathcal O_S$-\Term{module}{模}，$\mathcal F$ 为 $\mathcal A$ 的\Term{finite-type}{有限型} $\mathcal O_S$-\Term{submodule}{子模}，$\mathcal F^{**}$ 为其\Term{bidual}{双对偶}，$U$ 为 $\mathcal F$ 的\Term{flatness-locus}{平坦轨迹}，即使 $\mathcal F_s$ 为\Term{flat}{平坦} $\mathcal O_{S,s}$-\Term{module}{模}的点 $s\in S$ 所成的\Term{set}{集合}。

(i) $U$ 是 $S$ 的\Term{open-subset}{开集}，且 $\mathcal F^{**}=j_*j^*(\mathcal F)$，其中 $j$ 表示包含映射 $U\hookrightarrow X$。

(ii) 对任一\Term{quasi-coherent}{拟凝聚} $\mathcal O_U$-\Term{module}{模} $\mathcal E$，\Term{canonical-morphism}{典范态射}
\[
j_*(\mathcal E)\otimes_{\mathcal O_S}\mathcal A
\longrightarrow j_*(\mathcal E\otimes_{\mathcal O_U}j^*(\mathcal A))
\]
是\Term{isomorphism}{同构}。

(iii) 特别地，$\mathcal V=j_*j^*(\mathcal F)$ 是 $\mathcal A=j_*j^*(\mathcal A)$ 的\Term{submodule}{子模}，而\Term{canonical-morphism}{典范态射}
\[
\mathcal V\otimes_{\mathcal O_S}\mathcal A\longrightarrow j_*j^*(\mathcal F\otimes_{\mathcal O_S}\mathcal A)
\]
是\Term{isomorphism}{同构}。
\end{SGAresult}
\begin{proof}
用 $S$ 的一个\Term{connected-component}{连通分支}代替 $S$ 后，可假定 $S$ \Term{integral}{整}。根据 EGA IV$_2$ 2.1.12，$\mathcal F$ 的\Term{flatness-locus}{平坦轨迹}，即使 $\mathcal F_s$ 为\Term{flat}{平坦} $\mathcal O_{S,s}$-\Term{module}{模}的点 $s\in S$ 所成的\Term{set}{集合}，是 $S$ 的一个\Term{open-subset}{开集} $U$；记 $j$ 为包含映射 $U\hookrightarrow S$。由于 $\mathcal A_s$ \Term{flat}{平坦}从而\Term{torsion-free}{无挠}，$\mathcal F_s$ 也\Term{torsion-free}{无挠}，故 $U$ 包含全部\Term{codimension}{余维}不超过 $1$ 的点。因此，由 [BAC]，VII，§4.2，\Term{theorem}{定理} 1 的推论，有 $\mathcal F^{**}=j_*j^*(\mathcal F)$，从而得到\Term{monomorphism}{单态射} $\mathcal F^{**}\to j_*j^*(\mathcal A)$。

(ii) 的证明与 11.0 回顾的 EGA III 1.4.15 的证明类似。另一方面，由于 $S$ \Term{normal}{正规}，\Term{morphism}{态射} $\mathcal O_S\to j_*j^*(\mathcal O_S)$ 是\Term{isomorphism}{同构}（参见 EGA IV$_2$ 5.8.6 与 5.10.5）。将 (ii) 应用于 $\mathcal E=\mathcal O_U$，故有 $\mathcal A=j_*j^*(\mathcal A)$。最后，由于
\[
j^*(\mathcal F\otimes_{\mathcal O_S}\mathcal A)
=j^*(\mathcal F)\otimes_{\mathcal O_U}j^*(\mathcal A),
\]
(iii) 的最后一个断言由 (ii) 应用于 $\mathcal E=j^*(\mathcal F)$ 得出。\Term{lemma}{引理}得证。
\end{proof}

此外，回顾，一个\Term{finite-type}{有限生成} $R$-\Term{module}{模} $M$ 称为\Term{reflexive}{自反}，若从 $M$ 到其\Term{bidual}{双对偶} $M^{**}$ 的\Term{canonical-morphism}{典范态射}是\Term{isomorphism}{同构}。当 $R$ 是\Term{dimension}{维数}不超过 $2$ 的\Term{noetherian}{诺特}\Term{regular}{正则}\Term{ring}{环}时，这蕴含 $M$ 为\Term{projective-module}{投射模}。事实上，对任一\Term{finite-type}{有限生成} $R$-\Term{module}{模} $N$，考虑一个\Term{resolution}{分解} $L_1\to L_0\to N\to0$，其中 $L_0$ 与 $L_1$ 均为\Term{finite-free-module}{有限自由 $R$-模}；于是有\Term{exact-sequence}{正合列}
\[
0\longrightarrow N^*\longrightarrow L_0^*\longrightarrow L_1^*\longrightarrow Q\longrightarrow0,
\]
其中 $Q$ 表示 $L_0^*\to L_1^*$ 的\Term{cokernel}{余核}。又由于 $R$ 的\Term{homological-dimension}{同调维数}不超过 $2$（参见 [BAC]，X §4.2，\Term{theorem}{定理} 1 的推论 1），可得 $N^*$ 为\Term{projective-module}{投射模}。

\begin{SGAresult}{命题}{13.2}
设 $S$ 为\Term{dimension}{维数}不超过 $2$ 的\Term{noetherian}{诺特}\Term{regular}{正则}\Term{scheme}{概形}，$G$ 为\Term{affine}{仿射}\Term{flat}{平坦} $S$-\Term{group}{群}，$\mathscr A(G)$ 为其\Term{affine-algebra}{仿射代数}。
\SourcePage{440}
(i) 若 $G$ 在 $S$ 上为\Term{finite-type}{有限型}，则它\Term{isomorphism}{同构}于 $H=\underline{\operatorname{Aut}}_{\mathcal O_S}(\mathcal V)$ 的一个\Term{closed-subgroup}{闭子群}，其中 $\mathcal V$ 为某个\Term{finite-rank}{有限秩}\Term{locally-free}{局部自由} $\mathcal O_S$-\Term{module}{模}。若 $S$ 还为\Term{affine}{仿射}，则可取 $\mathcal V=\mathcal O_S^{\oplus d}$（某个 $d$），从而 $H=\operatorname{GL}_{d,S}$。

(ii) $\mathscr A(G)$ 是其各\Term{flat}{平坦}\Term{finite-type}{有限型}\Term{hopf-algebra}{Hopf 子 $\mathcal O_S$-代数}的\Term{filtered-inductive-limit}{滤过归纳极限}。
\end{SGAresult}
\begin{proof}
设 $\mathcal B$ 为 $\mathscr A(G)$ 的一个\Term{finite-type}{有限型}子 $\mathcal O_S$-\Term{algebra}{代数}。由于 $S$ 的一个\Term{open-subset}{开集}上的任一\Term{coherent}{凝聚}\Term{module}{模}都能延拓为 $S$ 上的\Term{coherent}{凝聚}\Term{module}{模}（参见 EGA I 9.4.5），存在 $\mathcal B$ 的一个\Term{coherent}{凝聚} $\mathcal O_S$-\Term{submodule}{子模} $\mathcal M$，它作为 $\mathcal O_S$-\Term{algebra}{代数}生成 $\mathcal B$（同上，9.6.5）。根据 11.10.bis，$\mathcal M$ 包含于一个 $G$-\Term{stable}{稳定}的\Term{coherent}{凝聚} $\mathcal O_S$-\Term{submodule}{子模} $\mathcal F$ 中。（注意，因为这里 $G$ 在 $S$ 上\Term{affine}{仿射}，该处的证明可写得更简单：可将其中的 $f_*f^*(\mathcal E)$ 换成 $\mathcal E\otimes_{\mathcal O_S}\mathscr A(G)$，等等。）

记 $j$ 为包含映射 $U\hookrightarrow S$，其中 $U$ 是 $\mathcal F$ 的\Term{flatness-locus}{平坦轨迹}。根据\Term{lemma}{引理} 13.1 及随后的回顾，$\mathcal V=j_*j^*(\mathcal F)$ 是 $\mathscr A(G)$ 的\Term{locally-free}{局部自由} $\mathcal O_S$-\Term{submodule}{子模}，且由于\Term{canonical-morphism}{典范态射}
\[
\mathcal V\otimes\mathscr A(G)\longrightarrow j_*j^*(\mathcal F\otimes\mathscr A(G))
\]
是\Term{isomorphism}{同构}，$\mathcal V$ 是 $\mathscr A(G)$ 的子 $G$-\Term{module}{模}。$G$ 在 $\mathcal V$ 上的\Term{action}{作用}于是诱导\Term{affine}{仿射} $S$-\Term{group-morphism}{群态射} $\rho_{\mathcal V}:G\to H=\underline{\operatorname{Aut}}_{\mathcal O_S}(\mathcal V)$，从而诱导\Term{hopf-algebra}{Hopf $\mathcal O_S$-代数}的\Term{morphism}{态射} $\phi_{\mathcal V}:\mathscr A(H)\to\mathscr A(G)$。记 $\mathcal A_{\mathcal V}$ 为 $\phi_{\mathcal V}$ 的像；它是 $H$ 的一个\Term{closed-subgroup}{闭子群} $G_{\mathcal V}$ 的\Term{affine-algebra}{仿射代数}，其中 $G_{\mathcal V}$ 就是 $\rho_{\mathcal V}$ 的\Term{closed-image}{闭像}。我们证明 $\mathcal A_{\mathcal V}$ 包含 $\mathcal B$。

问题在 $S$ 上是\Term{local-scheme}{局部}的，故可假定 $S=\operatorname{Spec}(R)$，且 $V=\Gamma(S,\mathcal V)$ 为以 $v_1,\ldots,v_n$ 为基的\Term{free-module}{自由 $R$-模}；此时 $H\simeq\operatorname{GL}_{n,R}$，$\mathscr A(H)$ 作为 $R$-\Term{algebra}{代数}由“\Term{matrix-coefficients}{矩阵系数}” $c_{ij}$ 及元素 $d^{-1}$ 生成，其中 $d$ 表示\Term{determinant}{行列式}。设 $\Delta$（相应地，$\varepsilon$）为 $A$ 的\Term{comultiplication}{余乘法}（相应地，\Term{augmentation}{增广}）。对 $j=1,\ldots,n$，写成
\[
\Delta(v_j)=\sum_{i=1}^n v_i\otimes a_{ij};
\]
则 $a_{ij}=\phi(c_{ij})$ 属于 $\operatorname{Im}(\phi)$。另一方面，由于 $V$ 是 $A$ 的 $R$-\Term{submodule}{子模}，可使用等式 $(\varepsilon\otimes\operatorname{id}_A)\circ\Delta=\operatorname{id}_A$，它蕴含 $v_j=\sum_i\varepsilon(v_i)a_{ij}$ 属于 $\operatorname{Im}(\phi)$。由于 $V$ 包含 $B=\Gamma(S,\mathcal B)$ 的一组\Term{generator}{生成元}，可得 $B\subset\operatorname{Im}(\phi)$。

若 $G$ 在 $S$ 上为\Term{finite-type}{有限型}，可取 $\mathcal B=\mathscr A(G)$，这时 $\phi$ 为\Term{surjective}{满射}，故 $S$-\Term{group-morphism}{群态射} $G\to H=\underline{\operatorname{Aut}}_{\mathcal O_S}(\mathcal V)$ 是\Term{closed-immersion}{闭浸入}。

若 $S$ 还为\Term{affine}{仿射}，则存在\Term{finite-rank}{有限秩}\Term{locally-free}{局部自由} $\mathcal O_S$-\Term{module}{模} $\mathcal V'$，使 $\mathcal V\oplus\mathcal V'=\mathcal O_S^d$ 作为 $\mathcal O_S$-\Term{module}{模}成立。把 $\mathcal V'$ 看作\Term{trivial-module}{平凡 $G$-模}，便可将 $\mathcal V$ 替换为 $\mathcal O_S^d$，从而得到 $G$ 是 $\operatorname{GL}_{d,S}$ 的一个\Term{closed-subgroup}{闭子群}。

最后回到任意\Term{affine}{仿射}\Term{flat}{平坦} $S$-\Term{group}{群} $G$ 的情形。根据 EGA I 9.4.9，$\mathscr A(G)$ 是其各\Term{coherent}{凝聚} $\mathcal O_S$-\Term{submodule}{子模} $\mathcal M$ 的并，因而也是上述各\Term{hopf-algebra}{Hopf 子 $\mathcal O_S$-代数} $\mathcal A_{\mathcal V}$ 的并，故得 (ii)。
\end{proof}

\begin{SGAresult}{例}{13.3}
设 $R$ 为\Term{discrete-valuation-ring}{离散赋值环}，其\Term{uniformizer}{一致化元}为 $\pi$，\Term{fraction-field}{分式域}为 $K$。考虑 $R$-\Term{group}{群}的\Term{filtered-projective-system}{滤过射影系}：
\[
\cdots\longrightarrow\mathbb G_{a,R}\xrightarrow{\times\pi}\mathbb G_{a,R}\xrightarrow{\times\pi}\mathbb G_{a,R}.
\]
\SourcePage{441}
它对应于\Term{inductive-system}{归纳系} $R[X_0]\to R[X_1]\to R[X_2]\to\cdots$，其中\Term{transition-morphism}{过渡态射}由 $X_n=\pi X_{n+1}$ 给出。其\Term{projective-limit}{射影极限} $G$ 是一个\Term{affine}{仿射}\Term{flat}{平坦}、非\Term{finite-type}{有限型}的 $R$-\Term{group-scheme}{群概形}；它的\Term{special-fiber}{特殊纤维}为平凡的，\Term{generic-fiber}{泛纤维}则为 $\mathbb G_{a,K}$；其\Term{affine-algebra}{仿射 $R$-代数} $\mathscr A(G)$ 是 $K[X]$ 中常数项系数属于 $R$ 的\Term{polynomial}{多项式}所成的\Term{subring}{子环}。（注意，在注记 11.10.1 中已遇到过此例。）注意 $G$ \Term{representable}{表示}如下\Term{functor}{函子}：它将任一 $R$-\Term{algebra}{代数} $B$ 对应到由 $B$ 中元素所成的序列 $(x_n)_{n\in\mathbb N}$ 的\Term{set}{集合}，其中对每个 $n$ 都有 $x_n=\pi x_{n+1}$。（特别地，每个 $x_n$ 都无限次可被 $\pi$ 除。）
\end{SGAresult}

现设 $S$ 为\Term{noetherian}{诺特}\Term{scheme}{概形}，使任一\Term{coherent}{凝聚} $\mathcal O_S$-\Term{module}{模}都是一个\Term{finite-type}{有限型}\Term{locally-free}{局部自由} $\mathcal O_S$-\Term{module}{模}的商；例如，当 $S$ 为\Term{separated}{分离}\Term{noetherian}{诺特}\Term{regular}{正则}\Term{scheme}{概形}时即如此，参见 SGA 6，第 II 讲 2.1.1 与 2.2.7。（可以证明，\Term{dimension}{维数}不超过 $1$ 的任一\Term{noetherian}{诺特}\Term{regular}{正则}\Term{scheme}{概形}也有此性质；但若 $S$ 是\Term{field}{域} $k$ 上原点加倍的\Term{affine-plane}{仿射平面}，此性质便不成立，参见同上，2.2.7.2。）

\begin{SGAresult}{定义}{13.4}
设 $G$ 为\Term{affine}{仿射}\Term{flat}{平坦} $S$-\Term{group}{群}。依照 R. W. Thomason（[Th87]，2.1），若对每个\Term{coherent}{凝聚} $G$-$\mathcal O_S$-\Term{module}{模} $\mathcal F$，都存在\Term{finite-rank}{有限秩}\Term{locally-free}{局部自由}的 $G$-$\mathcal O_S$-\Term{module}{模} $\mathcal E$ 及一个 $G$-\Term{equivariant}{等变}\Term{epimorphism}{满态射} $\mathcal E\to\mathcal F$，就称二元组 $(G,S)$ 具有\Term{equivariant-resolution-property}{等变分解性质}（propriété de résolution équivariante），或满足 (RE)。
\end{SGAresult}

在该文\Term{theorem}{定理} 3.1 中，Thomason 在 $G$ 在 $S$ 上\Term{essentially-free}{本质自由}的假设下证明了下述结果（参见下面的注记）。Gabber 向我们指出了下面这个更简单的证明，它不使用此假设。

\begin{SGAresult}{命题}{13.5}
设 $S$ 为\Term{noetherian}{诺特}\Term{scheme}{概形}，$G$ 为\Term{affine}{仿射}、\Term{flat}{平坦}且\Term{finite-type}{有限型}的 $S$-\Term{group}{群}，其\Term{affine-algebra}{仿射代数}为 $\mathscr A(G)$。假定 $(G,S)$ 满足 (RE)。则：

(i) $G$ \Term{isomorphism}{同构}于 $H=\underline{\operatorname{Aut}}_{\mathcal O_S}(\mathcal V)$ 的一个\Term{closed-subgroup}{闭子群}，其中 $\mathcal V$ 为某个\Term{finite-rank}{有限秩}\Term{locally-free}{局部自由} $\mathcal O_S$-\Term{module}{模}。

(ii) 若 $S$ 还为\Term{affine}{仿射}，可对某个 $n$ 取 $\mathcal V=\mathcal O_S^{\oplus d}$，从而 $H=\operatorname{GL}_{d,S}$。
\end{SGAresult}
\begin{proof}
证明与 13.2 类似。如该处一样，存在一个 $G$-\Term{stable}{稳定}的\Term{coherent}{凝聚} $\mathcal O_S$-\Term{submodule}{子模} $\mathcal F$，它作为 $\mathcal O_S$-\Term{algebra}{代数}生成 $\mathcal A=\mathscr A(G)$。用 $S$ 的一个\Term{connected-component}{连通分支}代替 $S$ 后，可假定 $S$ \Term{connected}{连通}。根据假设 (RE)，存在一个秩为 $n$ 的\Term{locally-free}{局部自由} $G$-$\mathcal O_S$-\Term{module}{模} $\mathcal E$ 及一个 $\mathcal A$-\Term{comodule}{余模}的\Term{epimorphism}{满态射} $\pi:\mathcal E\to\mathcal F$。置 $H=\underline{\operatorname{Aut}}_{\mathcal O_S}(\mathcal E)$，这是一个 $S$-\Term{group-scheme}{群概形}，\Term{local-scheme}{局部}\Term{isomorphism}{同构}于 $\operatorname{GL}_{n,S}$。$G$ 在 $\mathcal E$ 上的\Term{action}{作用}诱导\Term{affine}{仿射} $S$-\Term{group-morphism}{群态射} $\rho:G\to H$，它对应于\Term{hopf-algebra}{Hopf $\mathcal O_S$-代数}的\Term{morphism}{态射} $\phi:\mathscr A(H)\to\mathscr A(G)$。我们证明 $\rho$ 是\Term{closed-immersion}{闭浸入}。

问题在 $S$ 上是\Term{local-scheme}{局部}的，故可假定 $S=\operatorname{Spec}(R)$，且 $V=\Gamma(S,\mathcal E)$ 是以 $v_1,\ldots,v_n$ 为基的\Term{free-module}{自由 $R$-模}；此时 $H\simeq\operatorname{GL}_{n,R}$，$B=\Gamma(S,\mathscr A(H))$ 作为 $R$-\Term{algebra}{代数}由“\Term{matrix-coefficients}{矩阵系数}” $c_{ij}$ 及元素 $d^{-1}$ 生成，其中 $d$ 表示\Term{determinant}{行列式}。设 $\Delta$（相应地，$\varepsilon$）为
\SourcePage{442}
$A=\Gamma(S,\mathscr A(G))$ 的\Term{comultiplication}{余乘法}（相应地，\Term{augmentation}{增广}），$\mu:V\to V\otimes_R A$ 为 $V$ 上的 $A$-\Term{comodule}{余模结构}，并置 $F=\Gamma(S,\mathcal F)$。对 $j=1,\ldots,n$，写成
\[
\mu(v_j)=\sum_{i=1}^n v_i\otimes a_{ij},
\]
则 $\phi(c_{ij})=a_{ij}$。另一方面，由于 $\pi:V\to F$ 是 $A$-\Term{comodule}{余模}的\Term{morphism}{态射}，有
\[
\sum_{i=1}^n\pi(v_i)\otimes a_{ij}
=(\pi\otimes_R\operatorname{id}_A)(\mu(v_j))=\Delta(\pi(v_j)),
\]
故
\[
\pi(v_j)=(\varepsilon\otimes_R\operatorname{id}_A)\Delta(\pi(v_j))
=\sum_{i=1}^n\varepsilon\pi(v_i)\,a_{ij}
=\phi\left(\sum_{i=1}^n\varepsilon\pi(v_i)c_{ij}\right).
\]
所以 $\phi(B)$ 包含 $\pi(V)=F$，而 $F$ 作为 $R$-\Term{algebra}{代数}生成 $A$，故 $\phi$ 为\Term{surjective}{满射}。这表明 $\rho$ 是\Term{closed-immersion}{闭浸入}，从而得到断言 (i)。

若 $S$ 还为\Term{affine}{仿射}，则存在\Term{finite-rank}{有限秩}\Term{locally-free}{局部自由} $\mathcal O_S$-\Term{module}{模} $\mathcal V'$，使 $\mathcal V\oplus\mathcal V'=\mathcal O_S^d$ 作为 $\mathcal O_S$-\Term{module}{模}成立。把 $\mathcal V'$ 看作\Term{trivial-module}{平凡 $G$-模}，便可将 $\mathcal V$ 替换为 $\mathcal O_S^d$，从而得到 $G$ 是 $\operatorname{GL}_{d,S}$ 的一个\Term{closed-subgroup}{闭子群}。命题得证。
\end{proof}

\textbf{注记 13.6。}（a）为完整起见，沿用前述记号，简述 Thomason 的论证（[Th87]，\Term{theorem}{定理} 3.1）。$G$-\Term{equivariant}{等变}\Term{epimorphism}{满态射} $\pi:\mathcal E\to\mathcal F$ 诱导一个\Term{closed-immersion}{闭浸入} $\tau:G\to\mathbb V(\mathcal E)$，满足 $\tau(g'g)=\tau(g')\cdot g$（注意：$G$ 在 $\mathbb V(\mathcal E)$ 上\Term{right-action}{右作用}）；并且有\Term{isomorphism}{同构}
\[
H\simeq\underline{\operatorname{Aut}}_{\mathcal O_S}(\mathbb V(\mathcal E))^{\mathrm{op}},
\]
它与 $G$ 在 $\mathcal E$ 上的\Term{left-action}{左作用}和在 $\mathbb V(\mathcal E)$ 上的\Term{right-action}{右作用}相容。令 $N$ 为\Term{strict-transporter}{严格输送子}（transporteur strict）$\underline{\operatorname{Transp}}^{\mathrm{str}}_H(\tau(G),\tau(G))$；当 $G$ 在 $S$ 上\Term{essentially-free}{本质自由}时，由 6.2.4 e) 得 $N$ 是 $H$ 的\Term{closed-group-subscheme}{闭群子概形}，从而在 $S$ 上\Term{affine}{仿射}。此外，$\rho$ \Term{resolution}{分解}为一个 $S$-\Term{group-morphism}{群态射} $\rho':G\to N$。另一方面，对每个 $S'\to S$ 及 $h\in N(S')$，置 $\pi(h)=\tau(1)\cdot h$（其中 $1$ 为 $G(S')$ 的\Term{neutral-element}{单位元}）；这定义一个 $S$-\Term{scheme}{概形}的\Term{morphism}{态射} $\pi:N\to\tau(G)$，它是 $\rho'$ 的\Term{retraction}{收缩}（将 $G$ 等同于 $\tau(G)$）。由于 $N$ 在 $S$ 上\Term{separated}{分离}，可得 $\rho'$ 是\Term{closed-immersion}{闭浸入}。

（b）看来 [Th87]，\Term{theorem}{定理} 3.1 的证明需要假设 $G$ 在 $S$ 上\Term{essentially-free}{本质自由}，但原文中并无此假设（作者转而引用 $H$ \Term{essentially-free}{本质自由}这一事实）。不过，当 $G$ 为\Term{reductive}{约化}时，此假设成立（参见第 XXII 讲 5.7.8），所以该文推论 3.2 考虑的所有情形都满足此假设。特别地，Thomason 在该文 2.5 证明：若 $S$ \Term{separated}{分离}、\Term{noetherian}{诺特}、\Term{regular}{正则}且\Term{dimension}{维数}不超过 $2$，$G$ 为\Term{affine}{仿射}、\Term{finite-presentation}{有限表示}，而且 $\mathscr A(G)$ 是\Term{locally-projective}{局部投射} $\mathcal O_S$-\Term{module}{模}，则 $(G,S)$ 满足 (RE)；根据 13.5，这在一个稍强的假设下给出了 13.2。

最后指出，[Th87]，2.5 的证明可略作简化，如下所述。（为简便起见，置于 $S=\operatorname{Spec}(R)$ 为\Term{affine}{仿射}的情形。）

\SourcePage{443}
\begin{SGAresult}{命题}{13.7}
设 $R$ 为\Term{dimension}{维数}不超过 $2$ 的\Term{noetherian}{诺特}\Term{regular}{正则}\Term{ring}{环}，$C$ 为作为 $R$-\Term{module}{模}是\Term{projective-module}{投射模}的 $R$-\Term{coalgebra}{余代数}，$F$ 为在 $R$ 上\Term{finite-type}{有限生成}的 $C$-\Term{comodule}{余模}。则 $F$ 是一个在 $R$ 上\Term{finitely-generated-projective}{有限生成投射}的 $C$-\Term{comodule}{余模} $V$ 的商。
\end{SGAresult}
\begin{proof}
用 $\operatorname{Spec}(R)$ 的一个\Term{connected-component}{连通分支}代替它后，可假定 $R$ 为\Term{integral-domain}{整环}，其\Term{fraction-field}{分式域}为 $K$。用 $\Delta$（相应地，$\varepsilon$）表示 $C$ 的\Term{comultiplication}{余乘法}（相应地，\Term{augmentation}{增广}），用 $\rho:F\to F\otimes C$ 表示 $F$ 上的\Term{comodule}{余模结构}。设 $\pi:W\to F$ 为\Term{surjective}{满射}\Term{morphism}{态射}，其中 $W$ 为\Term{finite-rank}{有限秩}\Term{free-module}{自由 $R$-模}。赋予 $W\otimes C$ 由 $\operatorname{id}_W\otimes\Delta$ 定义的\Term{comodule}{余模结构}，对 $F\otimes C$ 亦如此。于是 $\rho:F\to F\otimes C$ 是 $C$-\Term{comodule}{余模}的\Term{morphism}{态射}，它以 $\operatorname{id}_F\otimes\varepsilon$ 为\Term{section}{截面}。

设 $W'$ 为下述\Term{cartesian-square}{笛卡儿方块}所定义的 $C$-\Term{comodule}{余模}：
\[
\begin{tikzcd}
W' \arrow[r] \arrow[d] & F \arrow[d,"\rho"]\\
W\otimes C \arrow[r,"\pi\otimes\operatorname{id}_C"'] & F\otimes C .
\end{tikzcd}
\]
即 $W'$ 等同于\Term{morphism}{态射} $W\otimes C\to(F\otimes C)/\rho(F)$ 的\Term{kernel}{核}，且由 $x\mapsto(\pi\otimes\varepsilon)(x)$ 给出的\Term{projection}{投影} $\pi':W'\to F$ 为\Term{surjective}{满射}。由于 $F$ 是\Term{finite-type}{有限生成} $R$-\Term{module}{模}，存在 $W'$ 的一个在 $R$ 上\Term{finite-type}{有限生成}的\Term{subcomodule}{子余模} $V'$，使 $\pi'(V')=F$。由于 $W\otimes C$ 无 $R$-\Term{torsion}{挠}，$V'$ 也无 $R$-\Term{torsion}{挠}；所以用 $V'$ 代替 $F$ 后，可在一开始就假定 $F$ \Term{torsion-free}{无挠}。

对这个新的 $F$ 应用前述构造，得到上述 $V'$。再考虑\Term{morphism}{态射}
\[
W\otimes C\longrightarrow E=\frac{W\otimes C\otimes K}{V'\otimes K}
\]
的\Term{kernel}{核}\Term{subcomodule}{子余模} $V$。

于是 $V$ 包含 $V'$，而 $Q=(W\otimes C)/V$ 是 $K$-\Term{vector-space}{向量空间} $E$ 的一个 $R$-\Term{submodule}{子模}；置 $Q'=E/Q$。由于 $W\otimes C$ 与 $E$ 都是\Term{flat}{平坦} $R$-\Term{module}{模}，对任一 $R$-\Term{module}{模} $N$ 有
\[
\operatorname{Tor}_1^R(V,N)\simeq\operatorname{Tor}_2^R(Q,N)\simeq\operatorname{Tor}_3^R(Q',N),
\]
而 $R$ \Term{regular}{正则}且\Term{dimension}{维数}不超过 $2$，故右端为零。这表明 $V$ 是\Term{flat}{平坦} $R$-\Term{module}{模}。最后证明 $V$ 为\Term{finite-type}{有限生成} $R$-\Term{module}{模}。置 $M=W\otimes C$；由定义，$V/V'$ \Term{isomorphism}{同构}于 $M/V'$ 的 $R$-\Term{torsion}{挠}\Term{submodule}{子模} $(M/V')_{\mathrm{tors}}$。

由于 $M$ 是\Term{projective-module}{投射 $R$-模}，存在一个\Term{projective-module}{投射 $R$-模} $P$，使 $M\oplus P$ 为\Term{free-module}{自由 $R$-模} $L$。于是 $(M/V')_{\mathrm{tors}}\simeq(L/V')_{\mathrm{tors}}$。另一方面，由于 $V'$ 为\Term{finite-type}{有限生成}，存在 $L$ 的一个\Term{direct-summand}{直和项} $L'\simeq R^n$，使 $V'\subset L'$，因而还有 $(L/V')_{\mathrm{tors}}\simeq(L'/V')_{\mathrm{tors}}$；后者是\Term{finite-type}{有限生成}的，因为 $L'/V'$ 是\Term{finite-type}{有限生成}的。因此 $V$ 是\Term{finite-type}{有限生成}\Term{flat}{平坦} $R$-\Term{module}{模}，故为\Term{finitely-generated-projective}{有限生成投射模}（因为 $R$ \Term{noetherian}{诺特}）。命题 13.7 得证。
\end{proof}

\SourcePage{444}
\SGAsection{}{参考文献}
\footnote{编者注（140）：本讲引用的新增参考文献。}

\noindent [An73] S. Anantharaman，《\Term{dimension}{维数} $1$ 的\Term{base}{基底}上的\Term{group-scheme}{群概形}、\Term{homogeneous-space}{齐性空间}与\Term{algebraic-space}{代数空间}》（\textit{Schémas en groupes, espaces homogènes et espaces algébriques sur une base de dimension 1}），Mém. Soc. Math. France \textbf{33}（1973），5--79。
\par\noindent [Ba55] I. Barsotti，《\Term{group-variety}{群簇}的一个结构\Term{theorem}{定理}》（\textit{Un teorema di struttura per le varietà gruppali}），Atti Acad. Naz. Lincei Rend. Cl. Sci. Fis. Mat. Nat. \textbf{18}（1955），43--50。
\par\noindent [BLR] S. Bosch、W. Lütkebohmert、M. Raynaud，《\Term{neron-model}{Néron 模型}》（\textit{Néron Models}），Springer-Verlag，1990。
\par\noindent [BAC] N. Bourbaki，《\Term{commutative-algebra}{交换代数}》（\textit{Algèbre commutative}），第 I--IV、V--VII 及 X 章，Masson，1985 及 1998。
\par\noindent [BTop] N. Bourbaki，《一般\Term{topology}{拓扑}学》（\textit{Topologie générale}），第 I--IV 章，Hermann，1971。
\par\noindent [Br09] M. Brion，《\Term{anti-affine}{反仿射}\Term{algebraic-group}{代数群}》（\textit{Anti-affine algebraic groups}），J. Algebra \textbf{321}（2009），第 3 期，934--952。
\par\noindent [Ch60] C. Chevalley，《\Term{algebraic-group}{代数群}的一个\Term{theorem}{定理}的证明》（\textit{Une démonstration d’un théorème sur les groupes algébriques}），J. Maths. Pure Appl. \textbf{39}（1960），307--317。
\par\noindent [Co02] B. Conrad，《Chevalley \Term{algebraic-group}{代数群}\Term{theorem}{定理}的一个现代证明》（\textit{A modern proof of Chevalley’s theorem on algebraic groups}），J. Ramanujan Math. Soc. \textbf{17}（2002），1--18。
\par\noindent [DG70] M. Demazure、P. Gabriel，《\Term{algebraic-group}{代数群}》（\textit{Groupes algébriques}），Masson \& North-Holland，1970。
\par\noindent [Gr73] L. Gruson，《\Term{noetherian-ring}{诺特环}上\Term{flat-module}{平坦模}的\Term{homological-dimension}{同调维数}》（\textit{Dimension homologique des modules plats sur un anneau noethérien}），Symposia Mathematica \textbf{XI}（1973），243--254。
\par\noindent [Kn71] D. Knutson，《\Term{algebraic-space}{代数空间}》（\textit{Algebraic spaces}），Lect. Notes Maths. \textbf{203}，Springer-Verlag，1971。
\par\noindent [Ma07] B. Margaux，《非阿贝尔 \Term{cech-cohomology}{Čech 上同调}中的\Term{limit}{极限}过程》（\textit{Passage to the Limit in Non-Abelian Čech Cohomology}），J. Lie Theory \textbf{17}（2007），591--596。
\par\noindent [MW94] A. Masuoka、D. Wigner，《\Term{hopf-algebra}{Hopf 代数}的\Term{faithfully-flat}{忠实平坦}性》（\textit{Faithful flatness of Hopf algebras}），J. Algebra \textbf{170}（1994），156--184。
\par\noindent [Per76] D. Perrin，《域上的\Term{quasi-compact}{拟紧}\Term{group-scheme}{群概形}的逼近》（\textit{Approximation des schémas en groupes, quasi-compacts sur un corps}），Bull. Soc. Math. France \textbf{104}（1976），323--335。
\par\noindent [Pes66] C. Peskine，《Zariski“主\Term{theorem}{定理}”的一个推广》（\textit{Une généralisation du « Main Theorem » de Zariski}），Bull. Sci. Math. \textbf{90}（1966），119--127。
\par\noindent [PY06] G. Prasad、J.-K. Yu，《关于拟\Term{reductive-group-scheme}{约化群概形}》（\textit{On quasi-reductive group schemes}），J. Alg. Geom. \textbf{15}（2006），507--549。
\par\noindent [Ray70a] M. Raynaud，《\Term{group-scheme}{群概形}与\Term{homogeneous-space}{齐性空间}上的\Term{ample-sheaf}{丰沛层}》（\textit{Faisceaux amples sur les schémas en groupes et les espaces homogènes}），Lect. Notes Maths. \textbf{119}，Springer-Verlag，1970。
\par\noindent [Ray70b] M. Raynaud，《\Term{henselian-local-ring}{Hensel 局部环}》（\textit{Anneaux locaux henséliens}），Lect. Notes Maths. \textbf{169}，Springer-Verlag，1970。
\par\noindent [RG71] M. Raynaud、L. Gruson，《\Term{flatness}{平坦性}与\Term{projective-object}{投射}性的判别准则》（\textit{Critères de platitude et de projectivité}），Invent. math. \textbf{13}（1971），1--89。
\par\noindent [Ro56] M. Rosenlicht，《\Term{algebraic-group}{代数群}的若干基本\Term{theorem}{定理}》（\textit{Some basic theorems on algebraic groups}），Amer. J. Math. \textbf{78}（1956），401--443。

\SourcePage{445}
\noindent [SS09] C. Sancho de Salas、F. Sancho de Salas，《\Term{principal-bundle}{主丛}、拟\Term{abelian-variety}{阿贝尔簇}与\Term{algebraic-group}{代数群}的结构》（\textit{Principal bundles, quasi-abelian varieties and structure of algebraic groups}），J. Algebra \textbf{322}（2009），第 8 期，2751--2772。
\par\noindent [Se68] J.-P. Serre，《\Term{split-reductive-group}{分裂约化群}\Term{scheme}{概形}的 Grothendieck 群》（\textit{Groupes de Grothendieck des schémas en groupes réductifs déployés}），Publ. math. I.H.É.S. \textbf{34}（1968），37--52。
\par\noindent [Se99] C. S. Seshadri，《Chevalley：一些回忆》（\textit{Chevalley : some reminiscences}），Transform. Groups \textbf{4}（1999），第 2--3 期，119--125。
\par\noindent [Ta72] M. Takeuchi，《Hopf \Term{ideal}{理想}与 \Term{Hopf-subalgebra}{Hopf 子代数}之间的一个对应》（\textit{A correspondence between Hopf ideals and sub-Hopf algebras}），Manuscripta Math. \textbf{7}（1972），251--270。
\par\noindent [Ta79] M. Takeuchi，《相对 Hopf 模——\Term{equivalence}{等价}与自由性判别准则》（\textit{Relative Hopf modules -- Equivalences and freeness criteria}），J. Algebra \textbf{60}（1979），452--471。
\par\noindent [Th87] R. W. Thomason，《任意\Term{base-scheme}{基概形}上的\Term{equivariant}{等变}\Term{resolution}{分解}、\Term{linear}{线性}化与 Hilbert 第十四问题》（\textit{Equivariant resolution, linearization, and Hilbert’s fourteenth problem over arbitrary base schemes}），Adv. Math. \textbf{65}（1987），16--34。

\SourcePage{446}
% 原 PDF 第446页为空白页。
